% ch4.tex — Laguerre Polynomials, Integral Operators and Applications
% Hand-typeset精修版，所有公式规范 LaTeX 化，图片嵌入。

\chapter{Laguerre Polynomials, Integral Operators and Applications}

\section{Introduction}

In this chapter, we will go deeper into the discussion on special polynomials
and on the relevant applications to physical and mathematical problems.
We will deal with the Laguerre family and use the relevant two-variable forms.
As already discussed for the case of Hermite polynomials, the introduction of
an extra variable in the Theory of classical Polynomials is particularly useful.
Even though they can be reduced to the standard form by means of a simple
transformation, a further variable is the key tool to develop a formulation
based on an operational formalism.

This chapter will consist essentially of two parts. In the first, we will
continue our discussion on ``popular'' orthogonal sets by introducing the
Laguerre polynomials (LP). In the second, we will discuss the Theory of
Special Polynomials by using a more general point of view involving the
Appell and Sheffer classification. We will however adopt a ``non orthodox''
formulation.

\section{Laguerre Polynomials: Definitions and Basic Properties}

We go back to the definition of the negative derivative operator (NDO),
which satisfies
\begin{equation}
\hat D_x^{-n}\,1 = \frac{x^n}{n!}.
\label{eq:4.1.1}
\end{equation}

The two-variable Laguerre polynomial is defined as
\begin{equation}
L_n(x,y) = \bigl(y - \hat D_x^{-1}\bigr)^n 1.
\label{eq:4.1.2}
\end{equation}

The use of the ordinary binomial expansion yields
\begin{equation}
L_n(x,y) = \sum_{r=0}^{n} \binom{n}{r} (-1)^r y^{\,n-r} \hat D_x^{-r}\,1
= \sum_{r=0}^{n} \frac{(-1)^r y^{\,n-r} x^r}{(n-r)!\,r!}.
\label{eq:4.1.3}
\end{equation}
For $y=1$ this reduces to the ordinary Laguerre polynomial $L_n(x)$.

\begin{exercise}
Prove the scaling identity
\begin{equation}
L_n(x,y) = y^n L_n\!\left(\frac{x}{y}\right).
\label{eq:4.1.4}
\end{equation}
\end{exercise}

\begin{figure}[h]
\centering
% 注：原书 Fig.4.1 为两组曲线图（y=1 与 y=-1/2），此处用文字占位
% 如有需要可补绘并替换为本地的 fig4_1a.pdf / fig4_1b.pdf
\setlength{\fboxsep}{6pt}\framebox[\textwidth]{\rule{0pt}{1.2cm}
  \textbf{Fig.\,4.1 (placeholder):} First 5 Laguerre polynomials vs $x$
  for (a) $y=1$ and (b) $y=-\tfrac12$.}
\caption{First 5 Laguerre polynomials vs $x$ for (a) $y=1$ and (b) $y=-\tfrac12$.}
\label{fig:4.1}
\end{figure}

\paragraph{Remark 1.} The NDO is an integral operator:
$\hat D_x^{-1} f(x) = \int_a^x f(x')\,dx'$.
If we do not specify the lower integration limit it means that $a=0$.
In Chapter~9 we will provide further comments on the NDO formalism.

\subsection*{Quasi-monomiality}

Following the definitions of ``quasi-monomiality'' introduced in Chapter~3,
the LP are quasi-monomials. The relevant multiplicative and derivative
operators are
\begin{equation}
\hat M = y - \hat D_x^{-1},
\qquad
\hat P = \hat{\mathcal L}\!\hat D_x = -\partial_x \hat D_x^{-1}.
\label{eq:4.1.7}
\end{equation}

The LP are therefore defined in terms of the multiplicative operator as
\begin{equation}
L_n(x,y) = \hat M^n 1.
\label{eq:4.1.8}
\end{equation}

We can also push the formalism and use the following definition
\begin{equation}
\hat{\mathcal L}\!\hat D_x
= -\partial_x \hat D_x^{-1},
\label{eq:4.1.9}
\end{equation}
where $\hat{\mathcal L}\!\hat D_x$ is sometimes called the \emph{Laguerre derivative}.
Its validity can be checked by noting that
\begin{equation}
\hat{\mathcal L}\!\hat D_x\,L_n(x,y)
= -\partial_x \hat D_x^{-1}(y-\hat D_x^{-1})^n 1
= n(y-\hat D_x^{-1})^{n-1}1
= n\,L_{n-1}(x,y).
\label{eq:4.1.10}
\end{equation}

From the definition of $\hat M$ and $\hat P$ we can infer the differential
equation satisfied by $L_n(x,y)$. Noting that $\hat D_x^{-1}\partial_x = \hat 1$, we find
\begin{equation}
\hat M\hat P\,L_n(x,y) = n\,L_n(x,y)
\label{eq:4.1.11}
\end{equation}
and then
\begin{equation}
y\,x\,L_n''(x,y) + (y-x)\,L_n'(x,y) = -n\,L_n(x,y).
\label{eq:4.1.12}
\end{equation}

The Laguerre derivative has the following remarkable property
\begin{equation}
\bigl(\hat{\mathcal L}\!\hat D_x\bigr)^n
= (-1)^n \partial_x^n x^n \partial_x^n,
\label{eq:4.1.13}
\end{equation}
which can be proved by induction.

\begin{exercise}
Use the previous properties to prove the inversion identity
\begin{equation}
\frac{x^n}{n!}
= \sum_{s=0}^{n} \binom{n}{s} (-1)^s y^s L_{n-s}(x,y).
\label{eq:4.1.14}
\end{equation}
\emph{Hint:} note that $\frac{x^n}{n!} = (y - (y-\hat D_x^{-1}))^n 1$.
\end{exercise}

According to eq.~\eqref{eq:4.1.13}, the repeated application of the
Laguerre derivative on the monomial $q_n(x) = (-1)^n \frac{x^n}{n!}$ yields
\begin{equation}
\bigl(\hat{\mathcal L}\!\hat D_x\bigr)^m q_n(x)
= \frac{n!}{(n-m)!}\,q_{n-m}(x), \qquad n>m.
\label{eq:4.1.15}
\end{equation}

\begin{exercise}
Check that the function
\begin{equation}
C_0(x) = \sum_{l=0}^{\infty} \frac{q_l(x)}{l!}
\label{eq:4.1.16}
\end{equation}
is an eigenfunction of the Laguerre derivative.
\end{exercise}

The function \eqref{eq:4.1.16} is called the \emph{0-th order Tricomi--Bessel
function}; its behaviour is shown in Fig.~\ref{fig:4.2}.

\begin{figure}[h]
\centering
\includegraphics[max width=0.7\linewidth,keepaspectratio]{fig4_2.png}
\caption{0-th order Tricomi--Bessel function $C_0(x)$.}
\label{fig:4.2}
\end{figure}

Before going further, the following correspondence should be kept in mind:
\begin{equation}
\partial_x \;\longleftrightarrow\; \hat{\mathcal L}\!\hat D_x,
\qquad
x^n \;\longleftrightarrow\; q_n(x),
\qquad
e^x \;\longleftrightarrow\; C_0(x).
\label{eq:4.1.17}
\end{equation}

The two-variable LP satisfy the partial differential equation
\begin{equation}
\begin{cases}
\partial_y L_n(x,y) = -\partial_x \hat D_x^{-1} L_n(x,y),\\[2pt]
L_n(x,0) = q_n(x),
\end{cases}
\label{eq:4.1.18}
\end{equation}
which can be easily stated using identity \eqref{eq:4.1.10}. This is the
analogue of the heat equation satisfied by Hermite polynomials.

From eq.~\eqref{eq:4.1.18} we can also define the LP through the
operational rule
\begin{equation}
L_n(x,y) = e^{-y\,\partial_x \hat D_x^{-1}}
\!\left[(-1)^n\frac{x^n}{n!}\right].
\label{eq:4.1.20}
\end{equation}

\section{Laguerre Polynomials Generating Function}

According to the definition of generating function, we write
\begin{equation}
G_L(x,y,t) = \sum_{n=0}^{\infty} \frac{t^n}{n!}
\bigl(y-\hat D_x^{-1}\bigr)^n 1
= e^{y t}\,e^{-t \hat D_x^{-1}} 1.
\label{eq:4.2.1}
\end{equation}

Since $y$ and $\hat D_x^{-1}$ commute, the exponential containing the
inverse derivative operator can be summed:
\begin{equation}
e^{-t \hat D_x^{-1}} 1
= \sum_{n=0}^{\infty} \frac{t^n}{n!}\,(-1)^n \hat D_x^{-n} 1
= \sum_{n=0}^{\infty} \frac{1}{n!}\,q_n(xt)
= C_0(xt).
\label{eq:4.2.2}
\end{equation}
Thus
\begin{equation}
G_L(x,y,t) = e^{yt}\,C_0(xt).
\label{eq:4.2.3}
\end{equation}

A further generating function can be derived using the multiplicative
operator:
\begin{equation}
S_L(x,y,t) = \sum_{n=0}^{\infty} t^n L_n(x,y)
= \frac{1}{1-yt}\,\exp\!\left(-\frac{xt}{1-yt}\right),
\qquad |yt|<1.
\label{eq:4.2.4}
\end{equation}

To prove this identity, from the operator definition we get
\begin{align}
S_L(x,y,t)
&= \sum_{n=0}^{\infty} (\hat M t)^n 1
 = \frac{1}{1-\hat M t}\,1
 = \frac{1}{1-yt}\,\frac{1}{1+\dfrac{\hat D_x^{-1}t}{1-yt}} \notag\\
&= \frac{1}{1-yt}\,
\sum_{r=0}^{\infty} \left(-\frac{\hat D_x^{-1}t}{1-yt}\right)^{\!r} 1
 = \frac{1}{1-yt}\,\exp\!\left(-\frac{xt}{1-yt}\right).
\label{eq:4.2.5}
\end{align}

The use of the above generating function and of eq.~\eqref{eq:4.1.20}
allows the derivation of the identity
\begin{equation}
e^{-y\,\partial_x \hat D_x^{-1}} e^{-xt}
= \frac{1}{1-yt}\,\exp\!\left(-\frac{xt}{1-yt}\right),
\label{eq:4.2.6}
\end{equation}
which is the analogue of the Glaisher identity (see eq.~\eqref{eq:2.2.18})
when exponential operators containing Laguerre derivatives are involved.

On the other side, the generating function $S_L$ in eq.~\eqref{eq:4.2.5}
allows us to establish the following integral representation for the two-variable
LP:
\begin{equation}
L_n(x,y) = \frac{e^{x/y}}{n!}\int_0^{\infty} e^{-s}\left(\frac{ys}{}\right)^{\!n}
C_0\!\left(\frac{xs}{y}\right) ds,
\label{eq:4.2.7}
\end{equation}
whose validity can be checked by using the generating function
\eqref{eq:4.2.4}.

\section{Orthogonality Properties of Laguerre Polynomials}

The Laguerre polynomials are specified, according to eq.~\eqref{eq:4.1.20},
by an operational definition analogous to that of the Hermite polynomials.
The following identity is a by-product of the same definition:
\begin{equation}
L_n(x,y+z) = e^{-z\,\partial_x \hat D_x^{-1}} L_n(x,y).
\label{eq:4.3.1}
\end{equation}

We can now prove that the LP provide an orthogonal set. We start from the
solution of the equation
\begin{equation}
\begin{cases}
\partial_y F(x,y) = -\partial_x \hat D_x^{-1} F(x,y),\\[2pt]
F(x,0) = g(x),
\end{cases}
\label{eq:4.3.2}
\end{equation}
which plays the same role as the heat equation in the case of Hermite
polynomials. Its solution can be written in the integral form
\begin{equation}
F(x,y) = \frac{e^{x/y}}{y}\int_0^{\infty} e^{-\sigma}\,
C_0\!\left(\frac{x\sigma}{y}\right) g(-y\sigma)\,d\sigma.
\label{eq:4.3.3}
\end{equation}

If $g(x) = \sum_{n=0}^{\infty} a_n \frac{x^n}{n!}$, we can write
\begin{equation}
F(x,y) = e^{-y\,\partial_x \hat D_x^{-1}}
\sum_{n=0}^{\infty} a_n \frac{x^n}{n!}
= \sum_{n=0}^{\infty} a_n (-1)^n L_n(x,y).
\label{eq:4.3.4}
\end{equation}

The use of identity \eqref{eq:4.2.7} allows us to cast the previous
equation as
\begin{equation}
F(x,y) = \sum_{n=0}^{\infty} (-1)^n a_n L_n(x,y)
= \frac{e^{x/y}}{y}\int_0^{\infty} e^{-\sigma}\,
C_0\!\left(\frac{x\sigma}{y}\right) g(-y\sigma)\,d\sigma,
\label{eq:4.3.5}
\end{equation}
which holds provided that summation and integration can be interchanged.

We can therefore proceed as in the previous section to get the set of
functions bi-orthogonal to the Laguerre family. For simplicity we consider
$y=1$ and assume that $g(x)$ can be expanded in series of LP:
\begin{equation}
g(x) = \sum_{n=0}^{\infty} \alpha_n L_n(x).
\label{eq:4.3.6}
\end{equation}

Applying $\partial_x \hat D_x^{-1}$ to the expansion yields
\begin{equation}
\partial_x \hat D_x^{-1} \Phi(x)
= \sum_{n=0}^{\infty} \frac{\alpha_n}{n!}\,(-x)^n.
\label{eq:4.3.7}
\end{equation}

The use of the Weierstrass transform yields
\begin{equation}
\partial_x \hat D_x^{-1} \Phi(x)
= e^{-x}\int_0^{\infty} e^{-\sigma}\,C_0(-x\sigma)\,g(\sigma)\,d\sigma.
\label{eq:4.3.8}
\end{equation}

These three equations finally give
\begin{equation}
\sum_{n=0}^{\infty} \frac{\alpha_n}{n!}\,(-x)^n
= e^{-x}\int_0^{\infty} e^{-\sigma}\,C_0(-x\sigma)\,g(\sigma)\,d\sigma.
\label{eq:4.3.9}
\end{equation}

Using the generating function of LP in terms of the 0-th order
Tricomi--Bessel function (see eq.~\eqref{eq:4.2.3}) allows us to write
\begin{equation}
\sum_{r=0}^{\infty} \frac{\alpha_r}{r!}\,(-x)^r
= \sum_{n=0}^{\infty} \frac{(-x)^n}{n!}
\int_0^{\infty} e^{-\sigma}\,L_n(\sigma)\,g(\sigma)\,d\sigma,
\label{eq:4.3.10}
\end{equation}
and thus
\begin{equation}
\alpha_n = \int_0^{\infty} e^{-\sigma}\,L_n(\sigma)\,g(\sigma)\,d\sigma.
\label{eq:4.3.11}
\end{equation}

At this point, it is not difficult to conclude that the bi-orthogonal
functions associated with the Laguerre polynomials are
\begin{equation}
\varphi_n(x) = e^{-x} L_n(x),
\label{eq:4.3.12}
\end{equation}
in such a way that
\begin{equation}
\int_0^{\infty} \varphi_n(x)\,L_m(x)\,dx = \delta_{nm}.
\label{eq:4.3.13}
\end{equation}

The functions $\varphi_n(x)$ satisfy a second-order differential equation
which can be directly inferred from that of the LP. We define
\begin{equation}
\hat h_{\mathrm{lag}}\,L_n(x) = 0,
\qquad
\hat h_{\mathrm{lag}} = x\partial_x^2 + (1-x)\partial_x + n,
\label{eq:4.3.14}
\end{equation}
and then the Hermitian associate
\begin{equation}
\hat h_{\mathrm{lag}}^\dagger \varphi_n(x) = 0,
\qquad
\hat h_{\mathrm{lag}}^\dagger = x\partial_x^2 + (1+x)\partial_x + n + 1.
\label{eq:4.3.15}
\end{equation}

The Laguerre orthogonal functions are
\begin{equation}
\ell_n(x) = e^{-x/2} L_n(x),
\label{eq:4.3.16}
\end{equation}
and play, in applications (Quantum Mechanics, Optics, \dots), the same role
as Hermite orthogonal functions. The relevant ODE can be written as
\begin{equation}
x\,\ell_n''(x) + \ell_n'(x)
+ \left(n - \frac{x}{4} + \frac{1}{2}\right)\ell_n(x) = 0.
\label{eq:4.3.17}
\end{equation}

\begin{figure}[h]
\centering
\includegraphics[max width=0.7\linewidth,keepaspectratio]{fig4_3.png}
\caption{The first 5 Laguerre orthogonal functions $\ell_n(x)$.}
\label{fig:4.3}
\end{figure}

In Fig.~\ref{fig:4.3} we have reported some examples of the behaviour of
the first few Laguerre orthogonal functions.

\section{Bessel Functions}

We have already introduced the function $C_0(x)$, which has been shown to
be an eigenfunction of the Laguerre derivative. It is the 0-th order of a
more general family of Bessel-type functions. By taking repeated derivatives
of $C_0(x)$ we find
\begin{equation}
\partial_x^n C_0(x) = (-1)^n C_n(x),
\label{eq:4.4.1}
\end{equation}
where
\begin{equation}
C_n(x) = \sum_{r=0}^{\infty} \frac{(-1)^r x^r}{r!\,(n+r)!}.
\label{eq:4.4.2}
\end{equation}
The functions $C_n(x)$ are higher-order Bessel-like functions of Tricomi type.

It is evident that the following recurrence follows from eq.~\eqref{eq:4.4.1}:
\begin{equation}
\partial_x C_n(x) = -C_{n+1}(x),
\label{eq:4.4.3}
\end{equation}
while from eq.~\eqref{eq:4.4.2} we get
\begin{equation}
n\,C_n(x) = C_{n-1}(x) - x\,\partial_x C_n(x)
= C_{n-1}(x) + x\,C_{n+1}(x).
\label{eq:4.4.4}
\end{equation}

From the above two recurrences we can derive the differential equation
satisfied by the Tricomi--Bessel functions. Eqs.~\eqref{eq:4.4.3}--\eqref{eq:4.4.4}
can be written as
\begin{equation}
\hat s_+ C_n(x) = C_{n+1}(x),
\qquad
\hat s_- C_n(x) = C_{n-1}(x),
\label{eq:4.4.5}
\end{equation}
where
\begin{equation}
\hat s_+ = -\partial_x,
\qquad
\hat s_- = \hat n + x\partial_x.
\label{eq:4.4.6}
\end{equation}
The operators $\hat s_\pm$ are called \emph{shift operators} and the number
operator $\hat n$ is defined by
\begin{equation}
\hat n\,C_m(x) = m\,C_m(x).
\label{eq:4.4.7}
\end{equation}

The ODE satisfied by $C_n(x)$ follows from
\begin{equation}
\hat s_+ \hat s_- C_n(x) = C_n(x),
\label{eq:4.4.8}
\end{equation}
which implies
\begin{equation}
x\,C_n''(x) + (1+n)\,C_n'(x) + C_n(x) = 0.
\label{eq:4.4.9}
\end{equation}

Tricomi--Bessel functions are also easily shown to be specified by the
generating function
\begin{equation}
B(x,t) = \sum_{n=-\infty}^{\infty} t^n C_n(x)
= e^{\,t - x/t}.
\label{eq:4.4.10}
\end{equation}

Furthermore, multiplying both sides of eq.~\eqref{eq:4.4.4} by $t^{n-1}$
we find
\begin{equation}
n\,t^{n-1} C_n = t^{n-1} C_{n-1} + \frac{x}{t^2}\,t^{n+1} C_{n+1},
\label{eq:4.4.11}
\end{equation}
from which we can derive the differential equation satisfied by $B(x,t)$:
\begin{equation}
\begin{cases}
\partial_t B(x,t) = \left(1 + \dfrac{x}{t^2}\right) B(x,t),\\[4pt]
B(x,0) = 1.
\end{cases}
\label{eq:4.4.12}
\end{equation}

Tricomi--Bessel functions are linked to the ordinary cylindrical Bessel
functions by the identity
\begin{equation}
C_n(x) = \left(\frac{x}{2}\right)^{-n} J_n\!\left(2\sqrt{x}\right),
\label{eq:4.4.13}
\end{equation}
which, once inverted, yields
\begin{equation}
J_n(x) = \left(\frac{x}{2}\right)^{\!n}
C_n\!\left(\frac{x^2}{4}\right).
\label{eq:4.4.14}
\end{equation}

The properties of the cylindrical Bessel functions can be deduced directly
from those of the Tricomi functions. They are summarized below.

\subsection*{Bessel $J_n(x)$: key properties}

\begin{enumerate}
\item \textbf{Series expansion}
\begin{equation}
J_n(x) = \sum_{r=0}^{\infty}
\frac{(-1)^r (x/2)^{n+2r}}{r!\,(n+r)!}.
\label{eq:4.4.15}
\end{equation}

\item \textbf{Reflection property}
\begin{equation}
J_{-n}(x) = J_n(-x) = (-1)^n J_n(x).
\label{eq:4.4.16}
\end{equation}

\item \textbf{Generating function}
\begin{equation}
J(x,t) = \sum_{n=-\infty}^{\infty} t^n J_n(x)
= \exp\!\left[\frac{x}{2}\left(t - \frac{1}{t}\right)\right].
\label{eq:4.4.17}
\end{equation}

\item \textbf{Recurrences}
\begin{align}
\partial_x J_n(x) &= \tfrac12\bigl(J_{n-1}(x) - J_{n+1}(x)\bigr),\\
n\,J_n(x) &= \tfrac{x}{2}\bigl(J_{n-1}(x) + J_{n+1}(x)\bigr).
\label{eq:4.4.18}
\end{align}

\item \textbf{Bessel differential equation}
\begin{equation}
x^2 J_n''(x) + x\,J_n'(x) + (x^2 - n^2)\,J_n(x) = 0.
\label{eq:4.4.19}
\end{equation}

\item \textbf{Jacobi generating function}
\begin{equation}
\sum_{n=-\infty}^{\infty} e^{in\vartheta} J_n(x) = e^{i x\sin\vartheta}.
\label{eq:4.4.20}
\end{equation}

\item \textbf{Integral representation}
\begin{equation}
J_n(x) = \frac{1}{\pi}\int_0^{\pi}
\cos\!\bigl(x\sin\vartheta - n\vartheta\bigr)\,d\vartheta.
\label{eq:4.4.21}
\end{equation}
\end{enumerate}

\begin{figure}[h]
\centering
\begin{minipage}{0.48\linewidth}
\includegraphics[max width=\linewidth,keepaspectratio]{fig4_4a.png}
\subcaption{$C_n(x)$ vs $x$.}
\end{minipage}\hfill
\begin{minipage}{0.48\linewidth}
\includegraphics[max width=\linewidth,keepaspectratio]{fig4_4b.png}
\subcaption{$C_n(x)$ vs $\log x$ (semi-log plot).}
\end{minipage}
\caption{Tricomi--Bessel functions $C_n(x)$: (a) the first 3 functions;
(b) log-linear plot showing the long-term oscillatory behaviour.}
\label{fig:4.4}
\end{figure}

\begin{figure}[h]
\centering
\setlength{\fboxsep}{6pt}\framebox[\textwidth]{\rule{0pt}{1.2cm}
  \textbf{Fig.\,4.5 (placeholder):} Comparison between the first three
  Tricomi--Bessel $C_n(x)$ and Bessel $J_n(x)$ functions.}
\caption{Comparison between the first three Tricomi--Bessel $C_n(x)$ and
Bessel $J_n(x)$ functions.}
\label{fig:4.5}
\end{figure}

\section{Associated Laguerre Polynomials}

In the previous sections we have treated the theory of Laguerre
polynomials and we have seen that their generating function is associated
with the 0-order Tricomi function. We may wonder whether new families
of LP can be obtained by modifying the generating function as follows:
\begin{equation}
\sum_{n=0}^{\infty} \frac{t^n}{(n+m)!}\,L_n^{(m)}(x,y)
= e^{yt}\,C_m(xt).
\label{eq:4.5.1}
\end{equation}

It can directly be checked that
\begin{equation}
L_n^{(m)}(x,y)
= \frac{(n+m)!}{n!}\,
\sum_{r=0}^{n} \frac{(-1)^r y^{\,n-r} x^r}{(n-r)!\,(m+r)!\,r!}.
\label{eq:4.5.2}
\end{equation}

The previous definition is also compatible with the operational form
\begin{equation}
L_n^{(m)}(x,y)
= (1 - y\partial_x)^m (y - \hat D_x^{-1})^n 1,
\label{eq:4.5.3}
\end{equation}
according to which we can derive most of the properties of this extension of
the LP, as for example the generating function
\begin{equation}
\sum_{n=0}^{\infty} t^n L_n^{(m)}(x,y)
= \frac{1}{(1-yt)^{m+1}}\,\exp\!\left(-\frac{xt}{1-yt}\right).
\label{eq:4.5.4}
\end{equation}

The Laguerre derivative for the associated Laguerre polynomials,
i.e. the operator satisfying
\begin{equation}
{}^{(m)}\!\hat{\mathcal L}\!\hat D_x\,
L_n^{(m)}(x,y) = (n+m)\,L_{n-1}^{(m)}(x,y),
\label{eq:4.5.5}
\end{equation}
is
\begin{equation}
{}^{(m)}\!\hat{\mathcal L}\!\hat D_x
= -\partial_x \hat D_x^{-1} - m\,\partial_x.
\label{eq:4.5.6}
\end{equation}

The relevant ODE can be derived in a fairly simple way starting from the
relation
\begin{equation}
L_n(x,y) = \frac{n!}{(n+m)!}\,\partial_x^m
\bigl[x^m L_n^{(m)}(x,y)\bigr].
\label{eq:4.5.7}
\end{equation}

Using the fact that the multiplicative operator has the same form as in the
ordinary ($m=0$) case yields
\begin{equation}
y\,x\,\partial_x^2 L_n^{(m)}
+ \bigl((1+m)y - x\bigr)\partial_x L_n^{(m)}
+ n\,L_n^{(m)} = 0.
\label{eq:4.5.8}
\end{equation}

We have now acquired sufficient skill to prove that the set of functions
($y=1$)
\begin{equation}
\ell_n^{(m)}(x) = \sqrt{\frac{n!}{(n+m)!}}\,e^{-x/2} x^{m/2}
L_n^{(m)}(x)
\label{eq:4.5.9}
\end{equation}
is orthogonal. The differential equation satisfied by these functions is
\begin{equation}
x\,\ell_n^{(m)\prime\prime}
+ \ell_n^{(m)\prime}
+ \left[n + \frac{1}{2} - \frac{x}{4} - \frac{m(m-2x)}{4x}\right]
\ell_n^{(m)} = 0.
\label{eq:4.5.10}
\end{equation}

Examples of the behaviour of the above functions for different values of
$m$ are shown in Fig.~\ref{fig:4.6}.

\begin{figure}[h]
\centering
\includegraphics[max width=0.7\linewidth,keepaspectratio]{fig4_6.png}
\caption{Laguerre orthogonal functions $\ell_n^{(m)}(x)$ for different values of
$m$ and $n$: (a) $n=3,\ m=0,4$; (b) $n=2,\ m=0,4$.}
\label{fig:4.6}
\end{figure}

\section{Legendre Polynomials}

We have claimed that the use of an extra variable in defining Laguerre or
Hermite functions is particularly useful since it yields a further degree of
freedom in the formalism. For Legendre polynomials $P_n(x)$ the
generating function is
\begin{equation}
\frac{1}{\sqrt{1-2xt+t^2}} = \sum_{n=0}^{\infty} t^n P_n(x),
\qquad |t|<1.
\label{eq:4.6.1}
\end{equation}

The Legendre polynomials satisfy the orthogonality relation
\begin{equation}
\int_{-1}^{1} P_n(x)\,P_m(x)\,dx = \frac{2}{2n+1}\,\delta_{nm},
\label{eq:4.6.2}
\end{equation}
and are solutions of Legendre's differential equation
\begin{equation}
(1-x^2)P_n''(x) - 2x\,P_n'(x) + n(n+1)\,P_n(x) = 0.
\label{eq:4.6.3}
\end{equation}

\section{Miscellaneous Applications and Comments}

Applications of Laguerre and Bessel functions include multipole expansions,
spherical Bessel functions in scattering theory, and the use of generating
functions to evaluate radial integrals in atomic physics. The associated
Laguerre functions appear naturally in the hydrogen-atom radial
wavefunctions, where the parameter $m$ is linked to the orbital angular
momentum quantum number $\ell$.

Further comments on Bessel functions and on their applications can be
found in Chapters~8--9.

\section{Appell Polynomials and Final Comments}

Appell polynomials $A_n(x)$ satisfy $A_n'(x) = n\,A_{n-1}(x)$ and form
the backbone of the umbral calculus treated in Chapter~9. The two-variable
Laguerre polynomials introduced in this chapter are indeed a particular
case of the Appell classification.

In Chapter~8 we will see how the same results can be achieved by the use
of a complementary technique, based on the ``umbral'' method.

\section*{Bibliography}

\begin{thebibliography}{99}

\bibitem{Lag1}
G. Szeg\H{o}, \emph{Orthogonal Polynomials}, 4th ed.,
Amer. Math. Soc. Colloq. Publ., Vol.~23, Providence, RI, 1975.

\bibitem{Lag2}
M. Abramowitz, I.A. Stegun (Eds.),
\emph{Handbook of Mathematical Functions},
Dover, New York, 1965.

\bibitem{Lag3}
W. Magnus, F. Oberhettinger, R.P. Soni,
\emph{Formulas and Theorems for the Special Functions of Mathematical Physics},
3rd ed., Springer, Berlin, 1966.

\bibitem{Lag4}
G. Dattoli, P.L. Ottaviani, A. Torre, L. Vazquez,
``Evolution operator equations, integration with algebraic and finite difference
methods: applications to physical problems in classical and quantum
mechanics,''
\emph{Riv. Nuovo Cimento} \textbf{20}, pp.~1--133, 1997.

\bibitem{Lag5}
G. Dattoli, S. Lorenzutta, A. Torre,
``Theory of multi-index multivariable Bessel functions and Hermite
polynomials,''
\emph{Matematiche} (Catania) \textbf{52}, pp.~179--197, 1997.

\bibitem{Lag6}
G. Dattoli, H.M. Srivastava, K. Zhukovsky,
``Orthogonality properties of the Hermite and related polynomials,''
\emph{J. Comput. Appl. Math.} \textbf{182}(1), pp.~165--172, 2005.

\bibitem{Lag7}
F.G. Tricomi,
\emph{Funzioni Ipergeometriche Confluenti},
Cremonese, Roma, 1954.

\bibitem{Lag8}
N.N. Lebedev,
\emph{Special Functions and Their Applications},
Dover Books on Mathematics, 1972.

\end{thebibliography}
